\section{课程结构、作业与项目}

\subsection{课程内容分层}

\begin{figure}[H]
\centering
\includegraphics[page=11,width=0.9\textwidth]{slides.pdf}
\caption{本课程把内容分成模型能力、Agent framework、应用和安全伦理四条主线。}
\end{figure}

课程结构分为四层：
\begin{enumerate}
    \item \textbf{Model core capabilities}：推理、规划、多模态理解。
    \item \textbf{LLM agent frameworks}：workflow、tool use、RAG、多 Agent。
    \item \textbf{Applications}：软件开发、工作流自动化、多模态和企业应用。
    \item \textbf{Safety \& ethics}：部署风险、人机交互、隐私和控制问题。
\end{enumerate}

\begin{knowledgebox}{项目方向的设计意图}
课程项目不是让学生重复做一个聊天机器人，而是鼓励围绕五类主题展开：应用、基准、基础能力、安全、多 Agent / 去中心化系统。这个设计本身就在提示学生，Agent 研究并不是单一方向，而是一个由能力、系统和治理共同组成的研究空间。
\end{knowledgebox}

\subsection{考核与项目组织}

导论中也给出了清晰的课程交付方式：
\begin{itemize}
    \item 每周阅读任务，用于建立跨讲次的共识背景；
    \item 一次 hands-on lab，把框架和工具真正跑起来；
    \item semester-long project，鼓励五人组队并围绕完整问题展开。
\end{itemize}

从质量要求上看，课程把 1--4 unit 分别对应到不同强度的实现深度，说明它更重视``把 Agent 系统做成一个可说明、可展示、可验证的工程实体''，而不只是阅读论文或写概念综述。

\subsection{本章小结}

这门课的组织方式与它的主题高度一致：课程本身就是一个从阅读、实验到项目落地的 Agent 系统训练场。

